\documentclass[../../../../main.tex]{subfiles}
\begin{document}

\begin{flushright}
\footnotesize\textit{【自動文字起こし・要確認】}
\end{flushright}

\noindent
[\textbf{解}] (i) $\frac{a}{\sqrt{m}} + \frac{b}{m}$ は $m$ について単調減少で ($\because m, a, b > 0$)
$m \to \infty$ で $0$ に収束し, $m=1$ で $a+b$ をとるから,

\bigskip
\noindent
$1^\circ$ $a+b < 1$ の時,
\[
4m < n^2 < 4m + \frac{a}{\sqrt{m}} + \frac{b}{m} < 4m+1
\]
となり, これが解答を持つとすると $4m$ と $4m+1$ の間に整数が存在することになり, 矛盾. よってこの時与不等式は解なし.

\bigskip
\noindent
$2^\circ$ $a+b \ge 1$ の時,
\[
\frac{a}{\sqrt{m_0}} + \frac{b}{m_0} \ge 1 > \frac{a}{\sqrt{m_0+1}} + \frac{b}{m_0+1}
\]
をみたす $m_0 \in \mathbb{N}$ が存在する. $m \ge m_0+1$ の時, $1^\circ$ と同じく与不等式は解なし. したがって解答があるのは $1 \le m \le m_0$ の時. さらにこの時,
\[
4m < N < 4m + \frac{a}{\sqrt{m}} + \frac{b}{m}
\]
をみたす $N \in \mathbb{N}$ は $\left[ \frac{a}{\sqrt{m}} + \frac{b}{m} \right]$ であるから, 全ての $N$ に $n \in \mathbb{Z}$ が 2つずつ対応していたとしても, $n$ は $2 \left[ \frac{a}{\sqrt{m}} + \frac{b}{m} \right]$ 個である. 以上から, 解は多くとも
\[
\sum_{m=1}^{m_0} 2 \left[ \frac{a}{\sqrt{m}} + \frac{b}{m} \right]
\]
であり, これは有限の値をとる.

\noindent
以上 $1^\circ, 2^\circ$ から示された. $\quad /\!/$

\bigskip
\noindent
(ロ) $4m < n^2 < 4m + \frac{8}{\sqrt{m}} + \frac{9}{m}$ ($m \ge 9$). $\quad \cdots \text{①}$

\noindent
$f(x) = \frac{8}{\sqrt{x}} + \frac{9}{x}$ とおく. $f'(x) = -\frac{4}{x\sqrt{x}} - \frac{9}{x^2} = -\frac{1}{x^2}(4\sqrt{x}+9)$ は $x \ge 9$ の時負の値をとるから, $x \ge 9$ で $f(x)$ は単調減少. これと $f(9) = \frac{11}{3}$ だから, $m \ge 9$ の時, $f(m) \le \frac{11}{3}$ である. 従って ①をみたす $n^2$ は
\[
n^2 = 4m+1, \, 4m+2, \, 4m+3
\]
に限られる. ここで, $\bmod 4$ で考えると, $n^2 \equiv 0, 1$ だから, 適するものは $n^2 = 4m+1$ に限られる.

\bigskip
\noindent
(ハ) $n>0$ の時を考えれば良く, この時 (ロ)から $n = \sqrt{4m+1}$, つまり $n$ が最大の時 $m$ が最大である. まず, $n^2 = 4m+1$ が解であるには
\[
1 < \frac{8}{\sqrt{m}} + \frac{9}{m}
\]
\[
\Leftrightarrow (m-9)^2 < 64m \quad (\because m \ge 9)
\]
が必要. これを解いて, $19 \le m \le 80$ とあわせて
\[
9 \le m \le 80 \quad \cdots \text{②}
\]
だから
\[
0 < 37 \le 4m+1 \le 321 \quad \cdots \text{③}
\]
であり, $36 < 37 \le 49$ から $6 < \sqrt{37} < 7$, $289 < 321 < 324$ から $17 < \sqrt{321} < 18$
なので, ③の両辺 $\sqrt{}$ をとって, $n \in \mathbb{N}$ から
\[
7 \le n \le 17
\]
従って $n=17$ をみたす $m$ があれば, これが求める値で, $n^2 = 4m+1$ を解いて計算すると $(n, m) = (17, 72)$ が求めるものである. $\quad /\!/$

\end{document}
